
\documentclass[a4paper,12pt]{article}

%%% Работа с русским языком
\usepackage{cmap}					% поиск в PDF
\usepackage{mathtext} 				% русские буквы в формулах
\usepackage[T2A]{fontenc}			% кодировка
\usepackage[utf8]{inputenc}			% кодировка исходного текста
\usepackage[english, russian]{babel}	% локализация и переносы
\usepackage{indentfirst}
\frenchspacing
\usepackage{amsmath} 
\usepackage{amsfonts}
\usepackage{amssymb}

\newcommand{\bz}{\mathbf{z}}
\newcommand{\bx}{\mathbf{x}}
\newcommand{\by}{\mathbf{y}}
\newcommand{\bv}{\mathbf{v}}
\newcommand{\bw}{\mathbf{w}}
\newcommand{\ba}{\mathbf{a}}
\newcommand{\bb}{\mathbf{b}}
\newcommand{\bp}{\mathbf{p}}
\newcommand{\bq}{\mathbf{q}}
\newcommand{\bt}{\mathbf{t}}
\newcommand{\bu}{\mathbf{u}}
\newcommand{\bT}{\mathbf{T}}
\newcommand{\bX}{\mathbf{X}}
\newcommand{\bZ}{\mathbf{Z}}
\newcommand{\bS}{\mathbf{S}}
\newcommand{\bH}{\mathbf{H}}
\newcommand{\bW}{\mathbf{W}}
\newcommand{\bY}{\mathbf{Y}}
\newcommand{\bU}{\mathbf{U}}
\newcommand{\bQ}{\mathbf{Q}}
\newcommand{\bP}{\mathbf{P}}
\newcommand{\bA}{\mathbf{A}}
\newcommand{\bB}{\mathbf{B}}
\newcommand{\bC}{\mathbf{C}}
\newcommand{\bE}{\mathbf{E}}
\newcommand{\bF}{\mathbf{F}}
\newcommand{\bomega}{\boldsymbol{\omega}}
\newcommand{\btheta}{\boldsymbol{\theta}}
\newcommand{\bgamma}{\boldsymbol{\gamma}}
\newcommand{\bdelta}{\boldsymbol{\delta}}
%\newcommand{\bPsi}{\boldsymbol{\Psi}}
%\newcommand{\bpsi}{\boldsymbol{\psi}}
\newcommand{\bxi}{\boldsymbol{\xi}}
\newcommand{\bchi}{\boldsymbol{\chi}}
\newcommand{\bzeta}{\boldsymbol{\zeta}}
\newcommand{\blambda}{\boldsymbol{\lambda}}
\newcommand{\beps}{\boldsymbol{\varepsilon}}
\newcommand{\bZeta}{\boldsymbol{Z}}
% mathcal
\newcommand{\cX}{\mathcal{X}}
\newcommand{\cY}{\mathcal{Y}}
\newcommand{\cW}{\mathcal{W}}

\newcommand{\dH}{\mathbb{H}}
\newcommand{\dR}{\mathbb{R}}
% transpose
\newcommand{\T}{^{\mathsf{T}}}

\renewcommand{\epsilon}{\ensuremath{\varepsilon}}
%\renewcommand{\phi}{\ensuremath{\varphi}}
\renewcommand{\kappa}{\ensuremath{\varkappa}}
\renewcommand{\le}{\ensuremath{\leqslant}}
\renewcommand{\leq}{\ensuremath{\leqslant}}
\renewcommand{\ge}{\ensuremath{\geqslant}}
\renewcommand{\geq}{\ensuremath{\geqslant}}
\renewcommand{\emptyset}{\varnothing}

%%% Дополнительная работа с математикой
\usepackage{amsmath,amsfonts,amssymb,amsthm,mathtools} % AMS
\usepackage{icomma} % "Умная" запятая: $0,2$ --- число, $0, 2$ --- перечисление

%% Номера формул
%\mathtoolsset{showonlyrefs=true} % Показывать номера только у тех формул, на которые есть \eqref{} в тексте.
%\usepackage{leqno} % Нумереация формул слева

%% Свои команды
\DeclareMathOperator{\sgn}{\mathop{sgn}}

%% Перенос знаков в формулах (по Львовскому)
\newcommand*{\hm}[1]{#1\nobreak\discretionary{}
	{\hbox{$\mathsurround=0pt #1$}}{}}

%%% Работа с картинками
\usepackage{graphicx}  % Для вставки рисунков
\setlength\fboxsep{3pt} % Отступ рамки \fbox{} от рисунка
\setlength\fboxrule{1pt} % Толщина линий рамки \fbox{}
\usepackage{wrapfig} % Обтекание рисунков текстом

%%% Работа с таблицами
\usepackage{array,tabularx,tabulary,booktabs} % Дополнительная работа с таблицами
\usepackage{longtable}  % Длинные таблицы
\usepackage{multirow} % Слияние строк в таблице

%%% Теоремы
\theoremstyle{plain} % Это стиль по умолчанию, его можно не переопределять.
\newtheorem{theorem}{Теорема}[section]
\newtheorem{proposition}[theorem]{Утверждение}

\theoremstyle{definition} % "Определение"
\newtheorem{corollary}{Следствие}[theorem]
\newtheorem{problem}{Задача}[section]

\theoremstyle{remark} % "Примечание"
\newtheorem*{nonum}{Решение}

%%% Программирование
\usepackage{etoolbox} % логические операторы

%%% Страница
\usepackage{extsizes} % Возможность сделать 14-й шрифт
\usepackage{geometry} % Простой способ задавать поля
\geometry{top=25mm}
\geometry{bottom=35mm}
\geometry{left=35mm}
\geometry{right=20mm}
%
%\usepackage{fancyhdr} % Колонтитулы
% 	\pagestyle{fancy}
%\renewcommand{\headrulewidth}{0pt}  % Толщина линейки, отчеркивающей верхний колонтитул
% 	\lfoot{Нижний левый}
% 	\rfoot{Нижний правый}
% 	\rhead{Верхний правый}
% 	\chead{Верхний в центре}
% 	\lhead{Верхний левый}
%	\cfoot{Нижний в центре} % По умолчанию здесь номер страницы

\usepackage{setspace} % Интерлиньяж
%\onehalfspacing % Интерлиньяж 1.5
%\doublespacing % Интерлиньяж 2
%\singlespacing % Интерлиньяж 1

\usepackage{lastpage} % Узнать, сколько всего страниц в документе.

\usepackage{soul} % Модификаторы начертания

\usepackage{hyperref}
\usepackage[usenames,dvipsnames,svgnames,table,rgb]{xcolor}
\hypersetup{				% Гиперссылки
	unicode=true,           % русские буквы в раздела PDF
	pdftitle={Заголовок},   % Заголовок
	pdfauthor={Автор},      % Автор
	pdfsubject={Тема},      % Тема
	pdfcreator={Создатель}, % Создатель
	pdfproducer={Производитель}, % Производитель
	pdfkeywords={keyword1} {key2} {key3}, % Ключевые слова
	colorlinks=true,       	% false: ссылки в рамках; true: цветные ссылки
	linkcolor=red,          % внутренние ссылки
	citecolor=black,        % на библиографию
	filecolor=magenta,      % на файлы
	urlcolor=cyan           % на URL
}

\usepackage{csquotes} % Еще инструменты для ссылок

%\usepackage[style=authoryear,maxcitenames=2,backend=biber,sorting=nty]{biblatex}

\usepackage{multicol} % Несколько колонок

\usepackage{tikz} % Работа с графикой
\usepackage{pgfplots}
\usepackage{pgfplotstable}


\begin{document}

\title{Использование интегрального преобразования в PINN}
\author{Боева Галина, группа М05-304а}
\date{\today}

\maketitle

\section{Введение}

Во многих вычислительных задачах в технике и науке дифференцирование функций или моделей имеет важное значение, но также необходимо и интегрирование. Важный класс вычислительных задач включает в себя так называемые интегро-дифференциальные уравнения, которые включают как интегралы, так и производные функции. В другом примере стохастические дифференциальные уравнения могут быть записаны в виде уравнения в частных производных для функции плотности вероятности стохастической переменной. Чтобы узнать характеристики для стохастической переменной, 
основанной на функции плотности, необходимо вычислить конкретные интегральные преобразования, а именно моменты, функции плотности. 
В последнее время парадигма машинного обучения в нейронных сетях, основанных на физике, приобретают все большую популярность как метод решения дифференциальных уравнений с использованием автоматического дифференцирования. В этой работе мы предлагаем дополнить парадигму нейронных сетей, основанных на физике, автоматической интеграцией, чтобы вычислять сложные интегральные преобразования для обученных решений и решать интегро-дифференциальные уравнения.

\section{Определение PINN}

Физически информированное машинное обучение предлагает решать дифференциальные уравнения, представляя приближенное решение с помощью искусственной нейронной сети, дифференцируемой квантовой схемы (DQC) или другого универсального функционального аппроксиматора (UFA) и оптимизируя UFA относительно функции потерь, построенной на основе дифференциальных уравнений. Параметризованное дифференциальное уравнение формально записывается как
\[
P(f, x; \lambda) = 0,
\]
где \(P\) — нелинейный оператор, действующий на \(f(x)\), а \(\lambda\) — параметр (обратите внимание, что \(x\) и \(\lambda\) могут быть многомерными). Если решение, аппроксимированное UFA, представлено как \(u_\theta\), то сходимость \(u_\theta\) к решению дифференциального уравнения (в большинстве случаев) обеспечивается условием
\[
\|P(u_\theta, x; \lambda)\| = 0.
\]
Функция потерь для решения дифференциальных уравнений в физически информированном машинном обучении вдохновлена вышеупомянутым наблюдением и определяется как
\[
\sum_{x_i} \|P(u_\theta, x_i; \lambda)\| = 0,
\]
где \(x_i\) представляет конечное число точек, выбранных из области определения дифференциального уравнения. UFA, такие как искусственные нейронные сети и дифференцируемые квантовые схемы (DQC), облегчают вычисление этой функции потерь с помощью автоматического дифференцирования.

\section{Автоматическая интегрирование}

Рассмотрим нейронную сеть
\[
f_\theta : \mathbb{R} \times \mathbb{R}^n \rightarrow \mathbb{R} \quad (s, v) \mapsto f_\theta(s, v).
\]
По основной теореме исчисления, частная антипроизводная \(F_\theta\) от \(f_\theta\) удовлетворяет
\[
f_\theta(s, v) = \int_{s_0}^s \frac{\partial f_\theta(q, v)}{\partial q} \, dq = \int_{s_0}^s F_\theta(s_0, v) \, ds_0.
\]
Обратите внимание, что \(F_\theta(s, v)\) также является функцией с той же областью и кодоменом, что и \(f_\theta(s, v)\), и может быть представлена нейронной сетью, которая использует те же параметры \(\theta\), что и \(f_\theta\). Lindell et al. отмечают, что можно оптимизировать \(F_\theta\) для представления некоторых целевых выходных данных и получить \(f_\theta\) без дополнительной оптимизации из оптимизированных параметров \(F_\theta\). 

Мы предлагаем дополнить парадигму PINN автоматическим интегрированием для вычисления интеграла функций, обученных с помощью PINN; кроме того, мы расширяем применение от интеграции функций до более общего понятия интегральных преобразований функций.

\section{Интегральные преобразования}

Интегральное преобразование — это математический оператор, который отображает функцию \(f\) в другую функцию \(\hat{f}\), определенную как
\[
\hat{f}(x, y) = \int k(x, x_0) f(x_0, y) \, dx_0,
\]
где функция \(k(x, x_0)\) известна как ядро интегрального преобразования.

Интегральное преобразование возникает в широком спектре областей. Оно часто используется для преобразования задачи, которая слишком сложна для решения в своем исходном представлении, например, для преобразования дифференциального уравнения в алгебраическое уравнение с помощью преобразования Фурье. Оно также часто возникает в областях, использующих вероятность или статистику, например, для гауссового сглаживания данных.

Если ядро является тождественным, интегральное преобразование сводится к обычному интегралу. Существуют различные стандартные преобразования, такие как преобразование Фурье и преобразование Лапласа, ядро которых фиксировано. Ядро фундаментальных решений дифференциальных уравнений и функций Грина, с другой стороны, может зависеть от дифференциальных уравнений и их граничных условий.

\section{Интегральные преобразования для функций, обученных с помощью PINN}
\begin{table}[h]
\centering
\begin{tabular}{|p{0.9\textwidth}|}
\hline
\textbf{Алгоритм} \\
\hline
1. Вход \\
Комбинация: (интегро-)дифференциальных уравнений, граничных условий, данных. \\
2. Преобразование \\
Преобразуйте весь вход с помощью: \\
$f(s, v) \rightarrow \frac{1}{k(s, s_0)} \frac{\partial G(q, s_0, v)}{\partial q} \bigg|_{q=s}$ \\
$\int_{s_b}^{s_a} k(s, s_0) f(s, v) \, ds \rightarrow G(s_b, s_0, v) - G(s_a, s_0, v)$. \\
3. Обучение \\
Обучите \(G\) с использованием дифференцируемой модели (PINN, DQC и т.д.). \\
4. Выход \\
Используйте обученную \(G\) для вычисления интегрального преобразования или неизвестной функции: \\
$\int_{s_b}^{s_a} k(s, s_0) f(s, v) \, ds = G(s_b, s_0, v) - G(s_a, s_0, v)$ \\
$f(s, v) = \frac{1}{k(s, s_0)} \frac{\partial G(q, s_0, v)}{\partial q} \bigg|_{q=s}$. \\
\hline
\end{tabular}
\caption{Алгоритм.}
\end{table}
В этой работе мы объединяем физически информированные (квантовые) нейронные сети и автоматическое интегрирование для формулировки квантово-совместимого метода машинного обучения для выполнения общего интегрального преобразования на обученной функции. Рассмотрим неизвестную функцию
\[
f : \mathbb{R} \times \mathbb{R}^n \rightarrow \mathbb{R} \quad (s, v) \mapsto f(s, v).
\]
Некоторая информация о функции \(f\) известна в виде:
\begin{itemize}
    \item значения \(f\) или ее производных в определенных частях области.
    \item интегро-дифференциальных уравнений неизвестной функции \(f\) в переменных \(s\) и \(v\).
\end{itemize}
Мы обсуждаем метод для представления в настройке PINN интегрального преобразования
\[
\hat{f}(s_0, v) = \int_{s_a(s_0)}^{s_b(s_0)} k(s, s_0) f(s, v) \, ds,
\]
где функции ядра \(k(s, s_0) \in \mathbb{R}\), \(s_a(s_0) \in \mathbb{R}\) и \(s_b(s_0) \in \mathbb{R}\) известны. Этот метод затем используется для:
\begin{itemize}
    \item оценки интегрального преобразования функции без явного обучения самой функции.
    \item решения интегро-дифференциальных уравнений для неизвестной функции \(f\).
\end{itemize}

Сначала мы определяем вспомогательную функцию
\[
G : \mathbb{R} \times \mathbb{R} \times \mathbb{R}^n \rightarrow \mathbb{R} \quad (s, s_0, v) \mapsto G(s, s_0, v).
\]

Если \(s_a\) и \(s_b\) являются функциями \(s_0\), мы определяем \(k(s, s_0) = 0\), если \(s < s_a(s_0)\) или \(s > s_b(s_0)\).

Следующим шагом является переформулировка данных и (интегро-)дифференциальных уравнений, заданных в терминах \(f\), в терминах вспомогательной функции \(G\). Это делается с помощью подстановок
\[
f(s, v) \rightarrow \frac{1}{k(s, s_0)} \frac{\partial G(q, s_0, v)}{\partial q} \bigg|_{q=s}
\]
\[
\int_{s_b}^{s_a} k(s, s_0) f(s, v) \, ds \rightarrow G(s_b, s_0, v) - G(s_a, s_0, v).
\]

С этими подстановками вся задача теперь сформулирована в терминах \(G\) и ее производных, без каких-либо ссылок на \(f\), ее производные или ее интеграл. Универсальный функциональный аппроксиматор (UFA) в настройке PINN, такой как классическая нейронная сеть или DQC, теперь может быть использован для прямого обучения \(G\). После завершения обучения интегральное преобразование неизвестной функции или сама неизвестная функция могут быть вычислены из уравнения (12). Сводка алгоритма показана в Таблице I.

\section{Интегральные преобразования на примере}

\subsection{Преобразование Фурье}

Преобразование Фурье функции \( f(x) \) определяется как:

\[
\mathcal{F}\{f(x)\} = \hat{f}(k) = \int_{-\infty}^{\infty} f(x) e^{-ikx} \, dx
\]

Обратное преобразование Фурье:

\[
\mathcal{F}^{-1}\{\hat{f}(k)\} = f(x) = \frac{1}{2\pi} \int_{-\infty}^{\infty} \hat{f}(k) e^{ikx} \, dk
\]

\subsection{Пример: Уравнение диффузии}

Рассмотрим уравнение диффузии:

\[
\frac{\partial u}{\partial t} = \nu \frac{\partial^2 u}{\partial x^2}
\]

Применим преобразование Фурье по переменной \( x \):

\[
\mathcal{F}\left\{ \frac{\partial u}{\partial t} \right\} = \frac{\partial \hat{u}(k, t)}{\partial t}
\]

\[
\mathcal{F}\left\{ \nu \frac{\partial^2 u}{\partial x^2} \right\} = -\nu k^2 \hat{u}(k, t)
\]

Таким образом, уравнение в частотном пространстве становится:

\[
\frac{\partial \hat{u}(k, t)}{\partial t} = -\nu k^2 \hat{u}(k, t)
\]

Это уравнение можно решить аналитически:

\[
\hat{u}(k, t) = \hat{u}(k, 0) e^{-\nu k^2 t}
\]

где \( \hat{u}(k, 0) \) — начальное условие в частотном пространстве.

\section{PINN и интегральные преобразования}

Предположим, у нас есть нейронная сеть \( \hat{u}(x, t; \theta) \), которая аппроксимирует \( u(x, t) \), где \( \theta \) — параметры сети.

Функция потерь \( L \) включает два компонента:

1. \textbf{Потеря на данных}:

\[
L_{\text{data}} = \frac{1}{N} \sum_{i=1}^N \left( \hat{u}(x_i, t_i; \theta) - u_i \right)^2
\]

2. \textbf{Потеря на физическом уравнении}:

\[
L_{\text{physics}} = \frac{1}{M} \sum_{j=1}^M \left( \frac{\partial \hat{u}(k_j, t_j; \theta)}{\partial t} + \nu k_j^2 \hat{u}(k_j, t_j; \theta) \right)^2
\]

Общая функция потерь:

\[
L = L_{\text{data}} + \lambda L_{\text{physics}}
\]

где \( \lambda \) — гиперпараметр, балансирующий влияние данных и физических законов.


\end{document}
